% ============================================================
% Topic 4 -- Number Sense -- part d: problems 41--50 (EN)
% Advanced extensions of problems 16--30
% ============================================================

\begin{probbox}{41}{\stars{1}}
\textbf{Sums of consecutive numbers} (extends Problem 16).
\begin{enumerate}[label=\alph*),leftmargin=1.6em,itemsep=1pt]
\item Find \emph{all} the ways of writing $100$ as a sum of at least two consecutive positive integers.
\item Prove that $64$ cannot be written as a sum of at least two consecutive positive integers.
\end{enumerate}
\end{probbox}
\probsol{
\textbf{Understanding the problem.} A sum of consecutive positive integers is determined by two pieces of information: \emph{how many numbers} (call it $k$, with $k\ge 2$) and \emph{the first number} (call it $a$, with $a\ge 1$). We need a quick formula for such a sum, and then we use it to ``narrow down'' the possibilities.

\begin{recallbox}\textbf{Sum of an arithmetic sequence} (Gauss's method): sum $=\dfrac{(\text{first}+\text{last})\times\text{number of terms}}{2}$.\end{recallbox}

\begin{proofbox}Write the sum twice, once forwards and once backwards, and add column by column: for example $S=9+10+\dots+16$ and $S=16+15+\dots+9$. Each column adds up to $9+16=25$ and there are $8$ columns, so $2S=25\cdot 8$, $S=100$. The same argument works for any arithmetic sequence.\end{proofbox}

\step{1}{Set up the key equation}
With $k$ consecutive numbers starting at $a$, the last number is $a+k-1$. Let $m=(\text{first}+\text{last})=2a+k-1$. By Gauss's formula: sum $=\dfrac{m\cdot k}{2}$. If the sum is $100$, then
\[k\cdot m=200.\]

\step{2}{Two important observations about $k$ and $m$}
\begin{itemize}[leftmargin=1.4em,itemsep=1pt]
\item $m-k=2a-1$ is odd, so \textbf{$k$ and $m$ have opposite parity} (one is even, the other odd).
\item $m-k=2a-1\ge 1$, so \textbf{$m>k$}.
\end{itemize}

\step{3}{List the factorisations $200=k\cdot m$ with $2\le k<m$}
Since $k<m$ and $k\cdot m=200$, we get $k\cdot k<200$, so $k\le 14$. The divisors of $200$ from $2$ to $14$ are $2,4,5,8,10$.
\begin{center}
\begin{tabular}{|c|c|c|c|}
\hline
\rowcolor{Tint}$k$ & $m=200\div k$ & Opposite parity? & Result\\
\hline
$2$ & $100$ & no (both even) & rejected\\
$4$ & $50$ & no & rejected\\
$5$ & $40$ & yes & $2a+4=40\Rightarrow a=18$\\
$8$ & $25$ & yes & $2a+7=25\Rightarrow a=9$\\
$10$ & $20$ & no & rejected\\
\hline
\end{tabular}
\end{center}

\step{4}{Write out the sums}
\[100=18+19+20+21+22\qquad\text{and}\qquad 100=9+10+11+12+13+14+15+16.\]
There are exactly \textbf{two} ways.

\begin{tipbox}Remember Problem 16: with $5$ consecutive numbers (an odd number of terms), the middle number equals the mean $=100\div 5=20$, so the five numbers are $18,19,20,21,22$. The ``middle number'' idea gives a very quick check.\end{tipbox}

\step{5}{Part b): why $64$ is impossible}
Suppose $64$ is a sum of $k\ge 2$ consecutive positive integers. Arguing as above: $k\cdot m=2\cdot 64=128$ with $m>k\ge 2$ and $k$, $m$ of opposite parity. So one of $k$, $m$ is odd and divides $128$. But $128=2^7$ has only \emph{one} odd divisor, namely $1$, while both $k$ and $m$ are $\ge 2$. Contradiction! So $64$ cannot be written in this way.

\begin{proofbox}The same argument shows that \emph{every power of $2$} ($2,4,8,16,32,64,\dots$) cannot be written as a sum of at least two consecutive positive integers. Conversely, every number with an odd divisor greater than $1$ can be.\end{proofbox}

\begin{warnbox}Do not forget the word ``positive'': if $0$ and negative numbers were allowed there would be extra ``fake'' ways, e.g.\ $(-17)+(-16)+\dots+21+22=100$ (the numbers from $-17$ to $17$ cancel, leaving $18+19+20+21+22$). Here every number must be $\ge 1$, so $a\ge 1$, which gives $m>k$.\end{warnbox}

\verify{$18+19+20+21+22=100$; $9+10+\dots+16=\dfrac{(9+16)\cdot 8}{2}=100$. Trying $64$ by hand: $2$ consecutive numbers have an odd sum; $3$ consecutive numbers have a sum divisible by $3$; \dots\ no way gives $64$.}
\ans{a) Two ways: $18+19+20+21+22$ and $9+10+\dots+16$. \ b) See the proof.}
}

\begin{probbox}{42}{\stars{1}}
\textbf{Common divisors and counting divisors} (extends Problem 17).
\begin{enumerate}[label=\alph*),leftmargin=1.6em,itemsep=1pt]
\item How many positive integers divide both $720$ and $1080$?
\item Find the smallest positive integer with exactly $12$ divisors.
\end{enumerate}
\end{probbox}
\probsol{
\textbf{Understanding the problem.} Listing the divisors of $720$ and $1080$ would be very long (each has dozens of divisors). We use two powerful tools: (1) the common divisors of two numbers are the divisors of their GCD; (2) the formula counting divisors from the prime factorisation.

\begin{recallbox}\textbf{Counting divisors.} If $n=p^{x}\cdot q^{y}\cdot r^{z}$ ($p,q,r$ different primes), then $n$ has $(x+1)(y+1)(z+1)$ divisors.\end{recallbox}

\begin{proofbox}Each divisor of $n$ has the form $p^{i}\cdot q^{j}\cdot r^{l}$ with $0\le i\le x$, $0\le j\le y$, $0\le l\le z$. The exponent $i$ has $x+1$ choices ($0,1,\dots,x$), $j$ has $y+1$ choices, $l$ has $z+1$ choices. By the multiplication rule there are $(x+1)(y+1)(z+1)$ divisors, and different choices give different divisors (because prime factorisation is unique).\end{proofbox}

\par\smallskip\textbf{Part a).}
\step{1}{Prime factorisation}
$720=2^4\cdot 3^2\cdot 5$ (since $720=16\cdot 45=16\cdot 9\cdot 5$); \quad $1080=2^3\cdot 3^3\cdot 5$ (since $1080=8\cdot 135=8\cdot 27\cdot 5$).

\step{2}{Find the GCD}
Take each common prime with the \emph{smaller} exponent: $\text{GCD}(720,1080)=2^3\cdot 3^2\cdot 5=8\cdot 9\cdot 5=360$.
(Another check: $1080-720=360$ and $720=2\cdot 360$, so the GCD is $360$.)

\step{3}{Count the divisors of $360$}
$360=2^3\cdot 3^2\cdot 5^1$ has $(3+1)(2+1)(1+1)=4\cdot 3\cdot 2=24$ divisors. So $24$ numbers divide both $720$ and $1080$.

\par\smallskip\textbf{Part b).}
\step{4}{The ``shapes'' of a number with exactly $12$ divisors}
The number of divisors is a product of factors (exponent $+1$). We write $12$ as a product of numbers $\ge 2$:
\[12=12,\qquad 12=6\cdot 2,\qquad 12=4\cdot 3,\qquad 12=3\cdot 2\cdot 2.\]
These correspond to four shapes: $p^{11}$, \ $p^{5}q$, \ $p^{3}q^{2}$, \ $p^{2}qr$.

\step{5}{The smallest number of each shape}
To make the number small, use the smallest primes $2,3,5$, and give the \emph{larger} exponent to the \emph{smaller} prime:
\begin{center}
\begin{tabular}{|c|c|c|}
\hline
\rowcolor{Tint}Shape & Smallest number & Value\\
\hline
$p^{11}$ & $2^{11}$ & $2048$\\
$p^{5}q$ & $2^5\cdot 3$ & $96$\\
$p^{3}q^{2}$ & $2^3\cdot 3^2$ & $72$\\
$p^{2}qr$ & $2^2\cdot 3\cdot 5$ & $\mathbf{60}$\\
\hline
\end{tabular}
\end{center}
The smallest is $60$.

\begin{warnbox}Do not give the large exponent to the large prime: for example $2\cdot 3^5=486$ is much bigger than $2^5\cdot 3=96$. And do not forget the shape with \emph{three} primes ($p^2qr$) -- it is the one that gives the answer.\end{warnbox}

\verify{The divisors of $60$: $1,2,3,4,5,6,10,12,15,20,30,60$ -- exactly $12$. A computer check of the numbers from $1$ to $59$: none has $12$ divisors.}
\ans{a) $24$ numbers. \ b) $60$.}
}

\begin{probbox}{43}{\stars{1}}
\textbf{How does a fraction change?} (extends Problems 18 and 19).
\begin{enumerate}[label=\alph*),leftmargin=1.6em,itemsep=1pt]
\item When the same whole number $n$ is added to both the numerator and the denominator of $\dfrac58$, the result is a fraction equal to $\dfrac45$. Find $n$.
\item Let $\dfrac ab$ be a fraction with $0<a<b$. Prove that for every positive integer $n$, $\dfrac{a+n}{b+n}>\dfrac ab$.
\item How many fractions equal to $\dfrac37$ have a two-digit denominator?
\end{enumerate}
\end{probbox}
\probsol{
\textbf{Understanding the problem.} Parts a) and b) are both about ``adding the same number to the numerator and the denominator''. What does \emph{not} change when we do that is the \emph{difference} between the denominator and the numerator -- this is the key.

\par\smallskip\textbf{Part a).}
\step{1}{Find the quantity that stays the same}
For $\dfrac58$, denominator $-$ numerator is $8-5=3$. After adding $n$ to both: $(8+n)-(5+n)=3$, still $3$.

\step{2}{Use a ``parts'' model as in Problem 19}
The new fraction equals $\dfrac45$, so the new numerator is $4$ parts and the new denominator is $5$ parts; the difference is $1$ part. Since the difference is $3$, $1$ part $=3$. New numerator $=4\cdot 3=12$, new denominator $=5\cdot 3=15$.

\step{3}{Find $n$}
$n=12-5=7$ (and also $15-8=7$).

\verify{$\dfrac{5+7}{8+7}=\dfrac{12}{15}=\dfrac{12\div3}{15\div3}=\dfrac45$. Correct.}

\par\smallskip\textbf{Part b).}
\step{4}{Compare ``what is missing to reach $1$''}
Since $a<b$, both fractions are less than $1$. What is missing to make $1$ is
\[1-\frac ab=\frac{b-a}{b},\qquad 1-\frac{a+n}{b+n}=\frac{(b+n)-(a+n)}{b+n}=\frac{b-a}{b+n}.\]
The fractions $\dfrac{b-a}{b}$ and $\dfrac{b-a}{b+n}$ have \emph{the same positive numerator}, and the denominator $b+n>b$. The fraction with the larger denominator is smaller, so $\dfrac{b-a}{b+n}<\dfrac{b-a}{b}$.

\step{5}{Conclusion}
The fraction $\dfrac{a+n}{b+n}$ is ``missing less'' to reach $1$, so it is larger: $\dfrac{a+n}{b+n}>\dfrac ab$.

\begin{recallbox}Two positive fractions with the same numerator: the one with the larger denominator is smaller (the same cake shared among more people gives smaller pieces).\end{recallbox}

\begin{proofbox}Another way (cross-multiplying): with positive denominators, $\dfrac{a+n}{b+n}>\dfrac ab$ exactly when $(a+n)\cdot b>a\cdot(b+n)$, i.e.\ $ab+nb>ab+na$, i.e.\ $nb>na$. This is true because $n>0$ and $b>a$.\end{proofbox}

\begin{warnbox}This is exactly why in Problem 18 we may \emph{not} ``add'' to the numerator and denominator: adding changes the fraction (a fraction less than $1$ gets bigger). Only \emph{multiplying} (or dividing) both by the same non-zero number keeps the fraction the same.\end{warnbox}

\par\smallskip\textbf{Part c).}
\step{6}{General form and condition}
Fractions equal to $\dfrac37$ have the form $\dfrac{3k}{7k}$ ($k$ a positive integer). The denominator $7k$ has two digits: $10\le 7k\le 99$. We have $7\cdot 1=7<10$, $7\cdot 2=14\ge 10$; $7\cdot 14=98\le 99$, $7\cdot 15=105>99$. So $k=2,3,\dots,14$: that is $14-2+1=13$ values.
The fractions are $\dfrac{6}{14},\dfrac{9}{21},\dots,\dfrac{42}{98}$.

\ans{a) $n=7$. \ b) See the proof. \ c) $13$ fractions.}
}

\begin{probbox}{44}{\stars{1}}
\textbf{Poles, a bicycle and a train} (extends Problem 20).
\begin{enumerate}[label=\alph*),leftmargin=1.6em,itemsep=1pt]
\item An cycles steadily at $18$ km/h along a road whose power poles are equally spaced $50$ m apart. How many minutes pass from the moment An passes the first pole to the moment he passes the $25$th pole?
\item A train travelling at constant speed takes $9$ seconds to pass a power pole and $30$ seconds to cross a bridge $350$ m long (from the moment the front of the train reaches the bridge until the end of the train has just left it). Find the length of the train and its speed in km/h.
\end{enumerate}
\end{probbox}
\probsol{
\textbf{Understanding the problem.} Part a) hides the ``fence-post trap'': $25$ poles make how many gaps? Part b) requires understanding how far a train (an object with length) must travel to ``pass'' something.

\begin{recallbox}Distance $=$ speed $\times$ time. Converting $1$ km/h: $1000$ m in $3600$ seconds, so $1$ km/h $=\dfrac{1000}{3600}=\dfrac{5}{18}$ m/s; conversely $1$ m/s $=3.6$ km/h.\end{recallbox}

\par\smallskip\textbf{Part a).}
\step{1}{Count the gaps between poles}
\begin{figC}
\foreach \x in {0,1,2,3}{\draw[Navy,line width=1.5pt] (\x*1.2,0)--(\x*1.2,0.8); \node[below,font=\scriptsize] at (\x*1.2,0) {pole \pgfmathparse{int(\x+1)}\pgfmathresult};}
\foreach \x in {0,1,2}{\draw[<->,tmOrange] (\x*1.2+0.05,0.4)--(\x*1.2+1.15,0.4); \node[above,font=\scriptsize,text=tmOrange!80!black] at (\x*1.2+0.6,0.4) {$50$ m};}
\node at (4.35,0.4) {$\cdots$};
\node[anchor=west,font=\footnotesize,text=Ink] at (5,0.4) {$4$ poles $\to$ $3$ gaps};
\end{figC}
Between the first and the $25$th pole there are $25-1=24$ gaps, with total length $24\cdot 50=1200$ m.

\step{2}{Convert the speed and find the time}
$18$ km/h $=18\cdot\dfrac{5}{18}=5$ m/s. Time: $1200\div 5=240$ seconds $=240\div 60=4$ minutes.

\begin{warnbox}Counting $25$ gaps ($25\cdot 50=1250$ m) gives $250$ seconds -- wrong. The number of gaps is always one less than the number of poles (counting from the first pole to the last).\end{warnbox}

\par\smallskip\textbf{Part b).}
\step{3}{What ``passing a pole'' and ``crossing a bridge'' mean}
\begin{itemize}[leftmargin=1.4em,itemsep=1pt]
\item Passing a pole: from when the front of the train is level with the pole until the end of the train is level with it, the train travels exactly \emph{one train length} $L$.
\item Crossing the bridge: from when the front is at the start of the bridge until the end leaves the far side, the train travels \emph{the bridge length plus the train length}: $350+L$.
\end{itemize}
\begin{figC}
\fill[Slate!25] (2,-0.15) rectangle (7,0); \node[below,font=\scriptsize] at (4.5,-0.15) {bridge $350$ m};
\draw[tmBlue,fill=tmBlue!20,rounded corners=2pt] (0,0.05) rectangle (2,0.45); \node[font=\scriptsize] at (1,0.25) {train (start)};
\draw[tmOrange,fill=tmOrange!20,rounded corners=2pt] (7,0.05) rectangle (9,0.45); \node[font=\scriptsize] at (8,0.25) {train (end)};
\draw[<->,Navy] (2,0.8)--(9,0.8) node[midway,above,font=\scriptsize] {the front travels $350+L$};
\draw[dashed,Slate] (2,0.05)--(2,0.8); \draw[dashed,Slate] (9,0.05)--(9,0.8);
\end{figC}

\step{4}{Compare the two situations}
Crossing the bridge takes $30$ seconds, passing the pole takes $9$ seconds. The difference $30-9=21$ seconds is exactly the time to travel the extra $350$ m (the bridge length). The speed of the train is
\[v=350\div 21=\frac{350}{21}=\frac{50}{3}\ \text{m/s}.\]

\step{5}{Find the train length and convert the speed}
Train length $L=v\cdot 9=\dfrac{50}{3}\cdot 9=150$ m. Speed: $\dfrac{50}{3}\cdot 3.6=\dfrac{50\cdot 3.6}{3}=60$ km/h.

\begin{warnbox}A common mistake: thinking the train travels only $350$ m to cross the bridge (forgetting to add the train length); then the speed would be $350\div 30\approx 11.7$ m/s, which is wrong. Follow \emph{one fixed point} on the train (e.g.\ its front) to know how far the train has travelled.\end{warnbox}

\verify{With $v=\dfrac{50}{3}$ m/s and $L=150$ m: passing the pole takes $150\div\dfrac{50}{3}=9$ seconds; crossing the bridge takes $(350+150)\div\dfrac{50}{3}=500\cdot\dfrac{3}{50}=30$ seconds. Both facts check out.}
\ans{a) $4$ minutes. \ b) The train is $150$ m long; its speed is $60$ km/h.}
}

\begin{probbox}{45}{\stars{2}}
\textbf{Products of digits} (extends Problem 21).
\begin{enumerate}[label=\alph*),leftmargin=1.6em,itemsep=1pt]
\item How many four-digit numbers have digits whose product is $36$?
\item Find the smallest positive integer whose digits have product $360$.
\item Find the largest positive integer with \emph{all different} digits whose product is $360$.
\end{enumerate}
\end{probbox}
\probsol{
\textbf{Understanding the problem.} All three parts start by \emph{factorising into primes}, then ``sharing'' the prime factors among the digits (each digit from $1$ to $9$). Note: no digit can be $0$ (otherwise the product is $0$).

\begin{recallbox}$36=2^2\cdot 3^2$; \ $360=2^3\cdot 3^2\cdot 5$. The digits and their factorisations: $1$; $2$; $3$; $4=2^2$; $5$; $6=2\cdot 3$; $7$; $8=2^3$; $9=3^2$.\end{recallbox}

\par\smallskip\textbf{Part a).}
\step{1}{Find the sets of four digits (ignoring order) with product $36$}
As in Problem 21, list $a\le b\le c\le d$:
\begin{itemize}[leftmargin=1.4em,itemsep=1pt]
\item Two digits equal to $1$: the other two digits have product $36$: $\{1,1,4,9\}$, $\{1,1,6,6\}$.
\item Exactly one digit $1$: the other three digits ($\ge 2$) have product $36$: $\{1,2,2,9\}$, $\{1,2,3,6\}$, $\{1,3,3,4\}$.
\item No digit $1$: four digits $\ge 2$ with product $36$: only $\{2,2,3,3\}$ (since $36=2\cdot2\cdot3\cdot3$ cannot be split further).
\item Three digits equal to $1$: the remaining digit would be $36$ -- impossible.
\end{itemize}

\begin{recallbox}\textbf{Arranging $4$ digits:} all four different: $4\cdot 3\cdot 2\cdot 1=24$ ways. Exactly two equal (the other two different from each other and from them): choose $2$ of the $4$ positions for the equal pair ($6$ ways), then place the other two digits in the $2$ remaining spots ($2$ ways): $12$ ways. Two equal pairs (like $1,1,6,6$): just choose $2$ positions for one pair: $6$ ways.\end{recallbox}

\step{2}{Count}
\begin{center}
\begin{tabular}{|c|c|c|}
\hline
\rowcolor{Tint}Digit set & Type & Arrangements\\
\hline
$\{1,1,4,9\}$ & one equal pair & $12$\\
$\{1,1,6,6\}$ & two equal pairs & $6$\\
$\{1,2,2,9\}$ & one equal pair & $12$\\
$\{1,2,3,6\}$ & all different & $24$\\
$\{1,3,3,4\}$ & one equal pair & $12$\\
$\{2,2,3,3\}$ & two equal pairs & $6$\\
\hline
\multicolumn{2}{|r|}{\textbf{Total}} & $\mathbf{72}$\\
\hline
\end{tabular}
\end{center}

\par\smallskip\textbf{Part b).}
\step{3}{Use as few digits as possible}
A number with fewer digits is smaller. One or two digits give a product of at most $9\cdot 9=81<360$, so we need at least $3$ digits.

\step{4}{Find three digits with product $360$}
$360$ has the prime factor $5$, and the only digit containing the factor $5$ is $5$ itself. So one digit is $5$, and the other two have product $360\div 5=72=8\cdot 9$ (the only way with digits $\le 9$). Digit set: $\{5,8,9\}$.

\step{5}{Arrange to get the smallest number}
Arrange the digits from smallest to largest: $589$.

\par\smallskip\textbf{Part c).}
\step{6}{For the largest number: as many digits as possible}
The digits are different, non-zero, and one of them must be $5$. Six different digits have product at least $1\cdot2\cdot3\cdot4\cdot5\cdot6=720>360$, so there are at most $5$ digits.

\step{7}{Find the sets of $5$ different digits with product $360$}
One digit is $5$; the other four are different, not $5$, with product $72=2^3\cdot 3^2$. Four different digits with a product as small as $72$ must include $1$ (since $2\cdot3\cdot4\cdot6=144>72$); the remaining three different digits have product $72$: $72=2\cdot4\cdot9=3\cdot4\cdot6$ (while $2\cdot 6\cdot 6$ and $3\cdot 3\cdot 8$ repeat a digit, and $1\cdot 8\cdot 9$ repeats the digit $1$). Two sets: $\{1,2,4,5,9\}$ and $\{1,3,4,5,6\}$.

\step{8}{Arrange to get the largest number}
Arrange from largest to smallest: $95421$ and $65431$. The larger is $95421$.

\begin{warnbox}In part c), do not jump to a number with a big first digit such as $985$ (product $360$): a $5$-digit number is always larger than a $3$-digit number. Maximise the \emph{number of digits} first, and only then the leading digit.\end{warnbox}

\verify{$5\cdot 8\cdot 9=360$; $9\cdot5\cdot4\cdot2\cdot1=360$. A computer search over all numbers with up to $6$ digits confirms: there are $72$ four-digit numbers with digit product $36$; the smallest number with digit product $360$ is $589$; the largest with all different digits and digit product $360$ is $95\,421$.}
\ans{a) $72$ numbers. \ b) $589$. \ c) $95\,421$.}
}

\begin{probbox}{46}{\stars{1}}
\textbf{Arithmetic grids} (extends Problem 22). In a square grid, the numbers in every row and in every column form arithmetic sequences.
\begin{enumerate}[label=\alph*),leftmargin=1.6em,itemsep=1pt]
\item A $4\times 4$ grid has these four corner numbers: top left $1$, top right $10$, bottom left $19$, bottom right $37$. Fill in the whole grid.
\item Find the sum of all $16$ numbers in the grid of part a). Compare it with the sum of the four corners: what do you notice? Explain.
\item Prove that in a $3\times 3$ grid of this kind, the number in the centre always equals the mean of the four corner numbers.
\end{enumerate}
\begin{figC}
\foreach \rr in {0,1,2,3}{\foreach \cc in {0,1,2,3}{\draw[Navy,line width=0.9pt,fill=white] (\cc*1.1,-\rr*0.8) rectangle (\cc*1.1+1.1,-\rr*0.8-0.8);}}
\node at (0.55,-0.4) {$1$}; \node at (3.85,-0.4) {$10$};
\node at (0.55,-2.8) {$19$}; \node at (3.85,-2.8) {$37$};
\end{figC}
\end{probbox}
\probsol{
\textbf{Understanding the problem.} Only the four corners are given. The first and last columns (and the first and last rows) each have a known first and last term, so they can be found immediately. Then we use the middle rows or columns.

\begin{recallbox}Arithmetic sequence with $4$ terms: difference $d=(\text{last}-\text{first})\div 3$. Sum of the $4$ terms $=2\times(\text{first}+\text{last})$ (pair term $1$ with term $4$, term $2$ with term $3$: both pairs have the same sum). With $3$ terms: middle term $=$ mean of the two outer terms.\end{recallbox}

\par\smallskip\textbf{Part a).}
\step{1}{First and last rows}
Row $1$: from $1$ to $10$, $d=(10-1)\div 3=3$: \ $1,4,7,10$.\\
Row $4$: from $19$ to $37$, $d=(37-19)\div 3=6$: \ $19,25,31,37$.

\step{2}{The four columns (first and last terms of each are known)}
Column $1$: $1\to 19$, $d=6$: $1,7,13,19$. \quad Column $2$: $4\to 25$, $d=7$: $4,11,18,25$.\\
Column $3$: $7\to 31$, $d=8$: $7,15,23,31$. \quad Column $4$: $10\to 37$, $d=9$: $10,19,28,37$.

\step{3}{Check the middle rows}
Row $2$: $7,11,15,19$ (difference $4$); row $3$: $13,18,23,28$ (difference $5$). The grid satisfies every condition:
\begin{figC}
\foreach \rr in {0,1,2,3}{\foreach \cc in {0,1,2,3}{\draw[Navy,line width=0.9pt,fill=white] (\cc*1.1,-\rr*0.8) rectangle (\cc*1.1+1.1,-\rr*0.8-0.8);}}
\foreach \cc/\val in {0/1,1/4,2/7,3/10}{\node at (\cc*1.1+0.55,-0.4) {$\val$};}
\foreach \cc/\val in {0/7,1/11,2/15,3/19}{\node at (\cc*1.1+0.55,-1.2) {$\val$};}
\foreach \cc/\val in {0/13,1/18,2/23,3/28}{\node at (\cc*1.1+0.55,-2.0) {$\val$};}
\foreach \cc/\val in {0/19,1/25,2/31,3/37}{\node at (\cc*1.1+0.55,-2.8) {$\val$};}
\foreach \rr/\val in {0/22,1/52,2/82,3/112}{\node[anchor=west,font=\footnotesize,text=Slate] at (4.55,-\rr*0.8-0.4) {row sum: $\val$};}
\end{figC}

\par\smallskip\textbf{Part b).}
\step{4}{Compute the sum}
Sum of the rows: $22+52+82+112=268$. Sum of the corners: $1+10+19+37=67$. We see that $268=4\cdot 67$: \textbf{the sum of the whole grid is $4$ times the sum of the corners.}

\step{5}{Explanation (true for every such $4\times4$ grid)}
Each row is an arithmetic sequence of $4$ terms, so each row sum $=2\times$(first term $+$ last term of the row). Adding the $4$ rows:
\[\text{Grid sum}=2\times(\text{sum of column }1+\text{sum of column }4).\]
Columns $1$ and $4$ are also arithmetic sequences of $4$ terms: sum of column $1=2\times(1+19)$, sum of column $4=2\times(10+37)$. So
\[\text{Grid sum}=2\times\big(2\cdot(1+19)+2\cdot(10+37)\big)=4\times(1+10+19+37)=4\times 67=268.\]

\par\smallskip\textbf{Part c).}
\step{6}{Name the cells of the $3\times 3$ grid}
\begin{figC}
\foreach \rr in {0,1,2}{\foreach \cc in {0,1,2}{\draw[Navy,line width=0.9pt,fill=white] (\cc*1.1,-\rr*0.8) rectangle (\cc*1.1+1.1,-\rr*0.8-0.8);}}
\node at (0.55,-0.4) {$a$}; \node at (1.65,-0.4) {$\cdot$}; \node at (2.75,-0.4) {$b$};
\node at (0.55,-1.2) {$u$}; \node[text=AnsRed] at (1.65,-1.2) {$x$}; \node at (2.75,-1.2) {$v$};
\node at (0.55,-2.0) {$c$}; \node at (1.65,-2.0) {$\cdot$}; \node at (2.75,-2.0) {$d$};
\end{figC}
The corners are $a,b,c,d$; the centre is $x$; $u$ and $v$ are the middle cells of the left and right columns.

\step{7}{Use ``middle term $=$ mean of the outer terms''}
The left column $a,u,c$ is arithmetic, so $u=\dfrac{a+c}{2}$. Right column: $v=\dfrac{b+d}{2}$. The middle row $u,x,v$ is arithmetic, so
\[x=\frac{u+v}{2}=\frac{\frac{a+c}{2}+\frac{b+d}{2}}{2}=\frac{a+b+c+d}{4}.\]
This is exactly the mean of the four corner numbers.

\begin{proofbox}Why does an arithmetic sequence $p,q,r$ satisfy $q=\dfrac{p+r}{2}$? Because $q-p=r-q$ (both equal the difference $d$), so $2q=p+r$.\end{proofbox}

\verify{Apply part c) to the top-left $3\times 3$ block of the grid in part a): corners $1,7,13,23$, mean $\dfrac{1+7+13+23}{4}=11$, exactly the centre $11$.}
\ans{a) See the grid. \ b) Sum $268=4\times 67$ ($4$ times the sum of the corners). \ c) See the proof.}
}

\begin{probbox}{47}{\stars{2}}
\textbf{``Rotating'' numbers} (extends Problem 23). In each part, every letter represents a digit, and $A\neq 0$.
\begin{enumerate}[label=\alph*),leftmargin=1.6em,itemsep=1pt]
\item Find the six-digit number $\nd{2ABCDE}$ such that $\nd{2ABCDE}\times 3=\nd{ABCDE2}$.
\item Find all digits $d$ (from $1$ to $9$) for which there is a six-digit number $\nd{dABCDE}$ with $\nd{dABCDE}\times 3=\nd{ABCDEd}$.
\item Prove that there is no six-digit number $\nd{1ABCDE}$ with $\nd{1ABCDE}\times 2=\nd{ABCDE1}$.
\end{enumerate}
\end{probbox}
\probsol{
\textbf{Understanding the problem.} All three parts are about ``moving the first digit to the end''. We reuse Method 1 of Problem 23: treat the block $\nd{ABCDE}$ as one number $N$ and set up an equation. This time we do it once for every leading digit $d$.

\begin{recallbox}With $N=\nd{ABCDE}$: \ $\nd{dABCDE}=d\cdot 100\,000+N$ \ and \ $\nd{ABCDEd}=10\cdot N+d$. Since $A\neq 0$, $N$ has exactly five digits: $10\,000\le N\le 99\,999$.\end{recallbox}

\step{1}{The general equation for parts a) and b)}
\[3\cdot(d\cdot 100\,000+N)=10N+d\ \Longrightarrow\ 300\,000\,d+3N=10N+d\ \Longrightarrow\ 299\,999\,d=7N.\]
Since $299\,999=7\cdot 42\,857$, dividing both sides by $7$ gives
\[N=42\,857\cdot d.\]

\step{2}{Part a): $d=2$}
$N=42\,857\cdot 2=85\,714$. The number is $285\,714$.

\step{3}{Part b): when does $N$ have five digits?}
$N=42\,857\,d\le 99\,999$. For $d=1$: $N=42\,857$ (five digits, fine). For $d=2$: $N=85\,714$ (fine). For $d=3$: $N=128\,571$ has six digits -- too big to be the block $\nd{ABCDE}$. For $d\ge 3$, $N$ is even bigger. So only $d=1$ and $d=2$ work, giving the two numbers $142\,857$ and $285\,714$.

\begin{tipbox}These two numbers use the same digits $1,4,2,8,5,7$ arranged in a circle. That is no accident: $\dfrac17=0.\overline{142857}$ and $\dfrac27=0.\overline{285714}$. Multiplying by $3$ corresponds to ``rotating'' the circle of digits.\end{tipbox}

\step{4}{Part c): set up the equation}
With $N=\nd{ABCDE}$:
\[2\cdot(100\,000+N)=10N+1\ \Longrightarrow\ 200\,000+2N=10N+1\ \Longrightarrow\ 199\,999=8N.\]

\step{5}{Part c): an odd/even contradiction}
The right-hand side $8N$ is even (since $8$ is even). The left-hand side $199\,999$ is odd. An odd number cannot equal an even number. So no such number exists.

\begin{proofbox}You can see it immediately without an equation: $\nd{1ABCDE}\times 2$ is even, so its last digit must be even; but $\nd{ABCDE1}$ ends in $1$ (odd). Contradiction.\end{proofbox}

\begin{warnbox}In part b), many students stop after finding $d=1$ (the answer of Problem 23) and forget $d=2$; or they accept $d=3,4,\dots$ without checking whether $N$ has more than five digits. Always check the condition ``$N$ has exactly five digits''.\end{warnbox}

\verify{$142\,857\times 3=428\,571$; \ $285\,714\times 3=857\,142$. In both cases the first digit has indeed moved to the end.}
\ans{a) $285\,714$. \ b) $d=1$ or $d=2$ (the numbers $142\,857$ and $285\,714$). \ c) See the proof.}
}

\begin{probbox}{48}{\stars{2}}
\textbf{A frog on the number line} (extends Problem 24). A frog sits at the point $0$ on the number line. On its $k$-th jump ($k=1,2,3,\dots$), the frog jumps exactly $k$ units, to the left or to the right as it likes.
\begin{enumerate}[label=\alph*),leftmargin=1.6em,itemsep=1pt]
\item After exactly $10$ jumps, can the frog be back at $0$? Why?
\item After exactly $10$ jumps, which points can the frog be at? How many such points are there?
\item What is the smallest number of jumps the frog needs to reach the point $20$?
\end{enumerate}
\end{probbox}
\probsol{
\textbf{Understanding the problem.} The frog's position after $n$ jumps is $\pm1\pm2\pm3\cdots\pm n$ (a $+$ sign for a jump to the right, a $-$ sign for a jump to the left). The idea is the same as in Problem 24: start from ``all plus'' and see how switching signs changes the result.

\begin{recallbox}$1+2+\dots+n=\dfrac{n(n+1)}{2}$ (Gauss's formula). For example $1+2+\dots+10=\dfrac{10\cdot 11}{2}=55$.\end{recallbox}

\step{1}{The ``switching'' rule}
If the frog always jumps right, after $10$ jumps it is at $55$. Switching one jump $k$ from right to left lowers the position by $2k$ (it loses $k$ and goes back another $k$). So if the jumps to the left add up to $s$, the final position is
\[55-2s.\]

\par\smallskip\textbf{Part a).}
\step{2}{Odd or even}
$55-2s$ is always \emph{odd} (odd minus even). The point $0$ is even, so the frog \textbf{cannot} be back at $0$ after $10$ jumps.

\par\smallskip\textbf{Part b).}
\step{3}{Which values can $s$ take?}
$s$ is the sum of a group (possibly empty) of different numbers from $\{1,2,\dots,10\}$, so $0\le s\le 55$. We show that \emph{every} integer $s$ from $0$ to $55$ can be made.

\begin{proofbox}Take the largest numbers one by one: $10, 9, 8, \dots$ The totals $0$, $10$, $10+9=19$, $19+8=27$, $27+7=34$, $34+6=40$, $40+5=45$, $45+4=49$, $49+3=52$, $52+2=54$, $54+1=55$ cut the interval $[0,55]$ into pieces. A number $s$ between two consecutive marks, e.g.\ $27<s\le 34$, is written as $s=10+9+8+(s-27)$ with $s-27$ from $1$ to $7$ -- a number not yet used. The same works for every other piece (for $0<s\le 10$ just take the number $s$). So every $s$ from $0$ to $55$ can be made.\end{proofbox}

\step{4}{The final positions}
For $s=0,1,2,\dots,55$, the position $55-2s$ takes the values $55,53,51,\dots,-53,-55$: that is, \textbf{every odd number from $-55$ to $55$}. Counting: from $-55$ to $55$ there are $28$ positive odd numbers ($1,3,\dots,55$) and $28$ negative ones, $56$ points altogether (the same as the number of values of $s$: $56$).

\begin{figC}
\draw[->,thick,Navy] (-5.9,0)--(5.9,0);
\foreach \x in {-5.5,-4.5,-3.5,-2.5,-1.5,-0.5,0.5,1.5,2.5,3.5,4.5,5.5}{\fill[KidOrange] (\x,0) circle (2.3pt);}
\draw[Navy] (0,0.1)--(0,-0.1); \node[below,font=\scriptsize] at (0,-0.1) {$0$};
\node[below,font=\scriptsize] at (-5.5,-0.1) {$-55$}; \node[below,font=\scriptsize] at (5.5,-0.1) {$55$};
\node[below,font=\scriptsize] at (0.5,-0.1) {$1$}; \node[below,font=\scriptsize] at (-0.5,-0.1) {$-1$};
\node[above,font=\scriptsize,text=Slate] at (0,0.15) {odd numbers only (sketch, not to scale)};
\end{figC}

\par\smallskip\textbf{Part c).}
\step{5}{Two necessary conditions}
After $n$ jumps, let $T=1+2+\dots+n$. The position is $T-2s$, so: (1) $T\ge 20$ (the frog cannot get farther than $T$); (2) $T-20$ must be even, i.e.\ $T$ must be even like $20$.

\step{6}{Try $n$ one at a time}
\begin{center}
\begin{tabular}{|c|c|c|c|c|c|c|c|}
\hline
\rowcolor{Tint}$n$ & $1$ & $2$ & $3$ & $4$ & $5$ & $6$ & $7$\\
\hline
$T$ & $1$ & $3$ & $6$ & $10$ & $15$ & $21$ & $28$\\
\hline
Reach $20$? & \multicolumn{5}{c|}{no ($T<20$)} & no ($T$ odd) & \textbf{possible}\\
\hline
\end{tabular}
\end{center}

\step{7}{Show a way with $n=7$}
We need $28-2s=20$, i.e.\ $s=4$: just jump left on the $4$th jump:
\[1+2+3-4+5+6+7=20.\]
So the smallest number of jumps is $7$.

\begin{warnbox}Checking only ``$T\ge 20$'' is not enough: for $n=6$, $T=21\ge 20$, but the position is always odd and can never be $20$. This is exactly the odd/even idea of Problem 24.\end{warnbox}

\verify{$1+2+3-4+5+6+7=20$. For $n=6$ the position is always $21-2s$, an odd number, never $20$. A computer search of all jump patterns also confirms: after $10$ jumps there are exactly $56$ positions, not including $0$; the smallest number of jumps to reach $20$ is $7$.}
\ans{a) No (the position is always odd). \ b) Every odd number from $-55$ to $55$: $56$ points. \ c) $7$ jumps.}
}

\begin{probbox}{49}{\stars{1}}
\textbf{Numbers containing the block $2025$} (extends Problem 28). Consider the six-digit numbers that contain the four digits $2,0,2,5$ \emph{together and in this order} (for example: $720\,251$, $202\,513$, $132\,025$). How many of these numbers are divisible by $45$?
\end{probbox}
\probsol{
\textbf{Understanding the problem.} The number consists of the block $2025$ plus two more digits $a,b$ in different positions. ``Divisible by $45$'' should be split into two conditions that are easier to check.

\begin{recallbox}$45=5\cdot 9$, and $5$ and $9$ have no common divisor other than $1$. So a number is divisible by $45$ exactly when it is divisible by both $5$ and $9$.
\textbf{Divisible by $5$:} the last digit is $0$ or $5$.
\textbf{Divisible by $9$:} the digit sum is divisible by $9$.\end{recallbox}

\step{1}{The condition for divisibility by $9$}
The digits of the block add up to $2+0+2+5=9$. The whole number has digit sum $9+a+b$, which is divisible by $9$ when $a+b$ is divisible by $9$, i.e.\ $a+b\in\{0,\ 9,\ 18\}$.

\step{2}{The three positions of the block, case by case}
\begin{itemize}[leftmargin=1.4em,itemsep=2pt]
\item \textbf{Form $\nd{2025ab}$} (the last digit is $b$, so $b\in\{0,5\}$; $a$ is any digit):
  $b=0$: $a\in\{0,9\}$ $\to$ $202\,500$, $202\,590$. \ $b=5$: $a+5\in\{9,18\}$ $\to$ $a=4$ $\to$ $202\,545$. \ \textbf{$3$ numbers.}
\item \textbf{Form $\nd{a2025b}$} ($a\neq 0$, $b\in\{0,5\}$):
  $b=0$: $a=9$ $\to$ $920\,250$. \ $b=5$: $a=4$ $\to$ $420\,255$. \ \textbf{$2$ numbers.}
\item \textbf{Form $\nd{ab2025}$} (ends in $5$ -- always divisible by $5$; $a\neq 0$):
  $a+b=9$: $a=1,2,\dots,9$ (with $b=9-a$) $\to$ $9$ numbers: $182\,025$, $272\,025$, \dots, $902\,025$. \ $a+b=18$: $a=b=9$ $\to$ $992\,025$. \ ($a+b=0$ is impossible because $a\neq0$.) \ \textbf{$10$ numbers.}
\end{itemize}

\step{3}{Check for double counting}
A number would be counted twice if the block $2025$ appeared twice among its six digits, i.e.\ two overlapping blocks. Shifting the block by $1$ place: ``$025$'' would have to match ``$202$'' -- false; shifting by $2$ places: ``$25$'' would have to match ``$20$'' -- false. So no number is counted twice.

\step{4}{Add up}
$3+2+10=15$ numbers.

\begin{warnbox}(1) Do not forget the case $a+b=0$ in the form $\nd{2025ab}$ (the number $202\,500$). (2) In the forms $\nd{a2025b}$ and $\nd{ab2025}$, the first digit $a$ cannot be $0$. (3) Check \emph{both} conditions, divisibility by $5$ and by $9$, not just one.\end{warnbox}

\verify{Example $992\,025$: it ends in $5$, and its digit sum $9+9+2+0+2+5=27$ is divisible by $9$, so it is divisible by $45$ ($992\,025=45\cdot 22\,045$). A computer search over all six-digit numbers also gives exactly $15$.}
\ans{$15$ numbers.}
}

\begin{probbox}{50}{\stars{2}}
\textbf{Mean and median} (extends Problem 30). A list of five positive integers (repeats allowed) has mean $12$ and median $10$.
\begin{enumerate}[label=\alph*),leftmargin=1.6em,itemsep=1pt]
\item What is the largest possible value of the largest number in the list?
\item What is the smallest possible value of the largest number in the list?
\item Answer both questions again if the five numbers must all be different.
\end{enumerate}
\end{probbox}
\probsol{
\textbf{Understanding the problem.} Arrange the five numbers in increasing order: $a\le b\le c\le d\le e$. The median is the middle number, so $c=10$. The mean is $12$, so the sum is $5\cdot 12=60$. We want to make $e$ as large or as small as possible. This is an ``optimisation'' problem: we must both \emph{prove} that a limit cannot be passed and \emph{give an example} that reaches exactly that limit.

\begin{recallbox}Mean $=$ sum $\div$ how many numbers, so sum $=12\cdot 5=60$. The median of five numbers (in increasing order) is the third one.\end{recallbox}

\step{1}{The key equation}
\[a+b+10+d+e=60\quad\Longrightarrow\quad a+b+d+e=50,\qquad a\le b\le 10\le d\le e.\]
\begin{barfig}
\foreach \x/\lab/\col in {0/a/tmBlue,1.3/b/tmBlue,2.6/10/tmGreen,3.9/d/tmOrange,5.2/e/tmOrange}{\blk{\x}{\x+1.2}{0}{0.6}{\col}{$\lab$}}
\node[anchor=west,font=\footnotesize] at (6.6,0) {sum $=60$, \ median $=10$};
\end{barfig}

\par\smallskip\textbf{Part a): make $e$ as large as possible.}
\step{2}{Upper limit}
For $e$ to be large, the other numbers must be as small as possible: $a\ge 1$, $b\ge 1$, $d\ge 10$. Hence $e=50-a-b-d\le 50-1-1-10=38$.
\step{3}{An example that reaches it}
The list $1,1,10,10,38$: sum $60$, median $10$. So the largest possible $e$ is $38$.

\par\smallskip\textbf{Part b): make $e$ as small as possible.}
\step{4}{Lower limit}
For $e$ to be small, $a$ and $b$ must be as large as possible, but $a\le b\le 10$, so $a+b\le 20$. Hence $d+e=50-(a+b)\ge 30$. Since $d\le e$, $2e\ge d+e\ge 30$, so $e\ge 15$.
\step{5}{An example that reaches it}
The list $10,10,10,15,15$: sum $60$, median $10$. So the smallest possible $e$ is $15$.

\par\smallskip\textbf{Part c): five different numbers} ($a<b<10<d<e$).
\step{6}{Largest $e$}
As small as possible: $a\ge 1$, $b\ge 2$, $d\ge 11$. So $e\le 50-1-2-11=36$. Example: $1,2,10,11,36$ (sum $60$). So the largest $e$ is $36$.
\step{7}{Smallest $e$}
As large as possible: $b\le 9$, $a\le 8$, so $a+b\le 17$ and $d+e\ge 50-17=33$. Since $d<e$, $d\le e-1$, so $2e-1\ge d+e\ge 33$, i.e.\ $e\ge 17$. Example: $8,9,10,16,17$ (sum $60$). So the smallest $e$ is $17$.

\begin{warnbox}A ``largest/smallest'' problem needs \emph{two} parts: (1) a proof that the limit cannot be passed, and (2) a concrete example reaching exactly that limit. An example alone may not be the best; a limit alone may not be reachable.\end{warnbox}

\begin{parentbox}Let your child share $60$ sweets among $5$ ``boxes'' kept in order, with the middle box always holding $10$. Ask: to make the last box as full as possible, should the other boxes hold few or many? And to make the last box as small as possible? This game helps children ``feel'' optimisation reasoning.\end{parentbox}

\verify{$1+1+10+10+38=60$; $10+10+10+15+15=60$; $1+2+10+11+36=60$; $8+9+10+16+17=60$; all have median $10$. A computer search over all lists also confirms the values $38$, $15$, $36$, $17$.}
\ans{a) $38$. \ b) $15$. \ c) Largest $36$, smallest $17$.}
}
